\textbf{Phần 1. Trắc nghiệm lựa chọn}
\setcounterpageref{ex}{0}

\begin{ex} %[3D16N76]
Số gồm $6$ chục nghìn, $4$ nghìn, $3$ trăm và $5$ đơn vị viết là
\choice
{$64350$}
{$64035$}
{$60435$}
{\True $64305$}
\loigiai{
Số đó là $64305$.\\
\textbf{Đáp án của câu hỏi là D}
}
\end{ex}

\begin{ex} %[3D16N76]
Số lớn nhất có 5 chữ số là
\choice
{$10000$}
{$100000$}
{\True $99999$}
{$99990$}
\loigiai{
Số lớn nhất có 5 chữ số là $99999$.\\
\textbf{Đáp án của câu hỏi là C}
}
\end{ex}

\begin{ex} %[3D16H76]
Chữ số La Mã $\text{XVIII}$ biểu thị số nào?
\choice
{$16$}
{\True $18$}
{$17$}
{$19$}
\loigiai{
$\text{XVIII} = 10 + 8 = 18$.\\
\textbf{Đáp án của câu hỏi là B}
}
\end{ex}

\begin{ex} %[3D16H76]
Làm tròn số $64\,850$ đến hàng nghìn ta được
\choice
{$64000$}
{$60000$}
{$64900$}
{\True $65000$}
\loigiai{
Hàng trăm là 8 nên làm tròn lên $65000$.\\
\textbf{Đáp án của câu hỏi là D}
}
\end{ex}

\begin{ex} %[3D16H76]
Làm tròn số $73\,400$ đến hàng chục nghìn ta được
\choice
{$80000$}
{$73000$}
{\True $70000$}
{$75000$}
\loigiai{
Hàng nghìn là 3 nên làm tròn thành $70000$.\\
\textbf{Đáp án của câu hỏi là C}
}
\end{ex}

\begin{ex} %[3D16H76]
Điền dấu thích hợp: $87\,450 \dots 87\,540$.
\choice
{$>$}
{$=$}
{Không so sánh được}
{\True $<$}
\loigiai{
Hàng trăm: $4 < 5$ nên $87450 < 87540$.\\
\textbf{Đáp án của câu hỏi là D}
}
\end{ex}

\begin{ex} %[3D16V76]
Số liền trước của số lớn nhất có 5 chữ số là
\choice
{\True $99998$}
{$99999$}
{$100000$}
{$99990$}
\loigiai{
$99999 - 1 = 99998$.\\
\textbf{Đáp án của câu hỏi là A}
}
\end{ex}

\begin{ex} %[3D16V76]
Từ 5 chữ số $9, 0, 4, 1, 7$ lập số lớn nhất có 5 chữ số khác nhau là
\choice
{$97401$}
{\True $97410$}
{$97140$}
{$94710$}
\loigiai{
Số lớn nhất là $97410$.\\
\textbf{Đáp án của câu hỏi là B}
}
\end{ex}

\begin{ex} %[3D16V76]
Tìm số tròn nghìn $x$ biết: $48\,500 < x < 50\,200$.
\choice
{\True $49000$ (hoặc $50000$)}
{$48000$}
{$51000$}
{$47000$}
\loigiai{
Số tròn nghìn thỏa mãn là $49000$ hoặc $50000$.\\
\textbf{Đáp án của câu hỏi là A}
}
\end{ex}

\begin{ex} %[3D16V76]
Số $45\,672$ gồm:
\choice
{$45$ nghìn, $6$ chục và $72$ đơn vị}
{$4$ nghìn, $5$ trăm, $6$ chục và $72$ đơn vị}
{\True $4$ chục nghìn, $5$ nghìn, $6$ trăm, $7$ chục và $2$ đơn vị}
{$456$ trăm và $72$ chục}
\loigiai{
$45672 = 40000 + 5000 + 600 + 70 + 2$.\\
\textbf{Đáp án của câu hỏi là C}
}
\end{ex}

\begin{ex} %[3D16V76]
Sắp xếp các số $32450; 34250; 32540; 34520$ theo thứ tự từ lớn đến bé:
\choice
{$32450 > 32540 > 34250 > 34520$}
{\True $34520 > 34250 > 32540 > 32450$}
{$34520 > 34250 > 32450 > 32540$}
{$34250 > 34520 > 32540 > 32450$}
\loigiai{
Thứ tự giảm dần: $34520 > 34250 > 32540 > 32450$.\\
\textbf{Đáp án của câu hỏi là B}
}
\end{ex}

\begin{ex} %[3D16C76]
Tìm số tròn chục nghìn nhỏ nhất có 5 chữ số mà tổng các chữ số của nó bằng 7.
\choice
{\True $70000$ (vì chữ số hàng chục nghìn bắt buộc phải là $7$)}
{$70000$}
{$16000$}
{$25000$}
\loigiai{
Số tròn chục nghìn có dạng $\overline{a0000} \Rightarrow a = 7 \Rightarrow 70000$.\\
\textbf{Đáp án của câu hỏi là A}
}
\end{ex}

\textbf{Phần 2. Trắc nghiệm đúng sai}
\setcounterpageref{ex}{0}

\begin{ex} %[3D16H76]
Xét tính đúng sai về số học:
\choiceTF
{\True Số $100\,000$ có 6 chữ số.}
{\True $\text{XIX}$ biểu thị số 19.}
{Làm tròn $35\,000$ đến hàng chục nghìn được $30\,000$.}
{\True Làm tròn $35\,000$ đến hàng chục nghìn được $40\,000$.}
\loigiai{
\begin{itemize}
    \item $100000$ có 6 chữ số: Đúng.
    \item $\text{XIX} = 19$: Đúng.
    \item Hàng nghìn là 5 nên làm tròn lên 40000. Ý 3 Sai, Ý 4 Đúng.
\end{itemize}
\begin{itemize}
    \item a) \textbf{Đúng}.
    \item b) \textbf{Đúng}.
    \item c) \textbf{Sai}.
    \item d) \textbf{Đúng}.
\end{itemize}
}
\end{ex}

\begin{ex} %[3D16V76]
Cho 3 số: $65\,400; 64\,500; 65\,040$. Xét tính đúng sai:
\choiceTF
{\True Số lớn nhất là $65\,400$.}
{\True Số bé nhất là $64\,500$.}
{Làm tròn số $65\,400$ đến hàng nghìn được $66\,000$.}
{\True Làm tròn số $65\,400$ đến hàng nghìn được $65\,000$ (vì hàng trăm là $4 < 5$).}
\loigiai{
\begin{itemize}
    \item Lớn nhất: $65400$. Đúng.
    \item Bé nhất: $64500$. Đúng.
    \item Hàng trăm là 4 nên làm tròn thành 65000. Ý 3 Sai, Ý 4 Đúng.
\end{itemize}
\begin{itemize}
    \item a) \textbf{Đúng}.
    \item b) \textbf{Đúng}.
    \item c) \textbf{Sai}.
    \item d) \textbf{Đúng}.
\end{itemize}
}
\end{ex}

\textbf{Phần 3. Trả lời ngắn}
\setcounterpageref{ex}{0}

\begin{ex} %[3D16H76]
Viết số liền trước của số $100\,000$.
\shortans{99999}
\loigiai{
\textbf{Đáp số:} $99999$\\
$100000 - 1 = 99999$.
}
\end{ex}

\begin{ex} %[3D16H76]
Chữ số La Mã $\text{XIV}$ có giá trị bằng bao nhiêu?
\shortans{14}
\loigiai{
\textbf{Đáp số:} $14$\\
14.
}
\end{ex}

\begin{ex} %[3D16V76]
Làm tròn số $76\,540$ đến hàng nghìn được số nào?
\shortans{77000}
\loigiai{
\textbf{Đáp số:} $77000$\\
$77000$.
}
\end{ex}

\begin{ex} %[3D16V76]
Làm tròn số $84\,900$ đến hàng chục nghìn được số nào?
\shortans{80000}
\loigiai{
\textbf{Đáp số:} $80000$\\
$80000$.
}
\end{ex}

\begin{ex} %[3D16V76]
Tìm số bé nhất có 5 chữ số khác nhau lập từ các chữ số $6, 0, 3, 8, 1$.
\shortans{10368}
\loigiai{
\textbf{Đáp số:} $10368$\\
$10368$.
}
\end{ex}

\begin{ex} %[3D16C76]
Tìm số tròn chục nghìn $x$ biết $55\,000 < x < 75\,000$.
\shortans{60000}
\loigiai{
\textbf{Đáp số:} $60000$\\
$60000$ (hoặc $70000$).
}
\end{ex}
